%% Errata to the submitted thesis, for the examiners. Page numbers refer to the submitted PDF
%% (submitted 30 August 2026; pagination checked against that PDF).
%% Build: pdflatex errata.tex
\documentclass[11pt]{article}
\usepackage[a4paper,margin=19mm]{geometry}
\usepackage[utf8]{inputenc}
\usepackage[T1]{fontenc}
\usepackage[british]{babel}
\usepackage{amsmath,amssymb}
\usepackage{booktabs,tabularx,array}
\usepackage{enumitem}
\usepackage[hidelinks]{hyperref}
\setlength{\parskip}{4pt}
\setlength{\parindent}{0pt}
\newcolumntype{L}{>{\raggedright\arraybackslash}X}
\newcolumntype{P}[1]{>{\raggedright\arraybackslash}p{#1}}

\begin{document}

{\large\bfseries Errata}\\[2pt]
\emph{Information Propagation Methods for Reliability, Reachability and Flow Analysis in
Directed Acyclic Process Networks}\\
Temitope Ohiani, Department of Civil and Environmental Engineering, University of Strathclyde\\
October 2026

\medskip
After submission, two corrections were made to the thesis. The first corrects a statement
about minimum cuts in Chapter~6; no number reported in the thesis changes. The second
corrects the capacity results of the Net3 case study in Chapter~10, and the statements that
depend on them in Chapter~11 and Appendix~B. Page numbers refer to the submitted copy, and a
copy of the corrected thesis with every change marked (deletions struck through in red,
additions underlined in blue) accompanies this sheet. The
reliability and schedule results of Chapter~10, and every result in the other chapters, are
unchanged. Every corrected value below was computed with version 0.2.3 of the registered
package, \texttt{InformationPropagationAnalysis.jl}, which contains the Chapter~6 correction.

\section*{1. Chapter 6: the structure of the minimum cuts}

Section~6.5.1 (pp.~105--106) stated that every node of the free zone $\mathcal{F}$ can be
assigned to either side of a minimum cut without changing its capacity, so that there are
exactly $2^{|\mathcal{F}|}$ minimum cuts. The converse is not true for every subset. By
Picard and Queyranne (1980), $S^\star \cup R$ with $R \subseteq \mathcal{F}$ is a minimum
cut if and only if $R$ is closed under reachability in the residual graph $G_f$. There are
therefore at most $2^{|\mathcal{F}|}$ minimum cuts, with equality only when no node of
$\mathcal{F}$ reaches another in $G_f$. For example, the power network of Appendix~A with
every capacity set to 60 has $\mathcal{F} = \{19, 20, 21\}$ linked as a chain in $G_f$, and
4 of its 8 subsets are minimum cuts.

\begin{tabularx}{\linewidth}{@{}P{3.4cm}L@{}}
\toprule
Location & Correction \\
\midrule
Def.~6.3 and 6.4, p.~105 & Taking the whole of $S^{\star\star}$ as the source side gives a
minimum cut; the claim that any free-zone node may be moved alone is removed. A theorem
stating the Picard--Queyranne characterisation is added. \\
p.~106, degeneracy & ``exactly $2^{|\mathcal{F}|}$ equivalent minimum cuts'' becomes ``at most
$2^{|\mathcal{F}|}$, one for each closed subset''. \\
p.~106, $E_{\text{some}}$ & The condition $v \notin \operatorname{Reach}_{G_f}(u)$ is added: a
saturated edge lies in some minimum cut only if its head is not reachable from its tail in
$G_f$. \\
p.~106, enumerator & The enumerator examines the subsets of $\mathcal{F}$ and keeps those whose
cut capacity equals $F^\star$; \texttt{is\_complete} refers to the subsets examined. \\
Table~6.2, p.~114 & Rows \texttt{edges\_some} ($O(|E_{\text{sat}}|(V+E))$, residual reach) and
\texttt{enumerate} (closed subsets of $\mathcal{F}$) updated. \\
\bottomrule
\end{tabularx}

The implementation was corrected to match, and the results of Chapter~6 were re-run. Configuration~A still has a free zone of two nodes and four minimum
cuts, of 13, 8, 8 and 3 crossing edges, and Configuration~B one. The IEEE RTS-24 free zone is
empty. No reported value changes.

\section*{2. Chapter 10: capacity results of the Net3 case study}

Section~10.3.2 specifies a super-sink with one edge from each of the 58 demand junctions,
with capacity equal to that junction's demand, 680.1 litres per second in total. The
capacity analysis reported in the thesis was run on the edge list without the super-sink and
its 58 edges. The reported throughput, 1,837.9 litres per second, therefore exceeded the total
demand. The analysis algorithms of the package were not at fault: they computed the maximum
flow of the graph they were given. The error lay in the interface layer. The case study's
edge list omitted the super-sink, and the server through which the interface reaches the
package (Section~8.8, p.~162) ignored the 58 capacity entries for edges it did not contain,
without reporting them. The
analysis was re-run on the edge list with the super-sink added and declared as the only sink,
and the results were checked against an independent max-flow computation (Edmonds--Karp,
separate implementation), which agrees to the reported precision. The server now rejects a
capacity file that names an edge the edge list does not contain.

\begin{tabularx}{\linewidth}{@{}P{3.6cm}LL@{}}
\toprule
Quantity & Submitted & Corrected \\
\midrule
Baseline throughput & 1,837.9 L/s & 680.1 L/s, the whole demand \\
Baseline minimum cut & unique; four edges: two discharges at the river pump's junction, one
at the lake's, and the source tank outlet & unique (free zone empty); the 58 demand edges \\
Saturated edges; paths & 7; 158 & 60; 401 \\
Sensitivity, baseline & the four cut edges, 883.3 to 76.0 L/s; edge 95--96 marginal range
unbounded & the demand edges, ranked by demand: junction 58, 280.1 L/s; junction 92, 103.3
L/s; every marginal range finite \\
Degraded throughput & 954.6 L/s, a 48\% drop & 43.3 L/s, 6.4\% of demand \\
Degraded minimum cuts & the same four edges & 24 cuts, free zone of six nodes; in every cut,
the river pump (95--97) and the demand edges of junctions 7, 8, 77, 79, 80, 81 \\
Flow timing, Table~B.5 & 2.33 s & 4.07 s \\
\bottomrule
\end{tabularx}

\medskip
\textbf{Interpretation.} At baseline the network meets its whole demand, so the network does
not bound delivery; the submitted text attributed the bound to the pumps and the source tank
(p.~189). In the degraded scenario only the six junctions downstream of the lake pump and the
discharging tank in the oriented graph can be served, and the derated trunk main does not
bind, because the river pump is its only feed. This result holds with every pipe in its
snapshot direction; a degraded state would need an orientation of its own, and this is
stated as a limitation of the operating-snapshot assumption.

\textbf{Profile overlay.} Section~10.7 (pp.~192--195) stated that the minimum cut and the
critical path have no node in common. The corrected figure draws the edges in every minimum
cut of the degraded scenario against the baseline critical path. They share nodes 95 and 97,
the ends of the river pump, and the super-sink. The pump whose loss leaves 6.4 per cent of
demand served lies on the restoration critical chain. The submitted overlay also shared nodes
95 and 97; the interface counted the overlap of an edge set by node only, and this is
corrected.

\begin{tabularx}{\linewidth}{@{}P{3.6cm}L@{}}
\toprule
Location & Change \\
\midrule
\S10.1, p.~182 & ``uploaded once'' removed; the capacity analysis reads the edge list with the
super-sink added. \\
\S10.3, pp.~183--184 & The super-sink is added to the capacity edge list (98 nodes, 177 edges)
and named as the only sink. \\
\S10.5, pp.~189--191; Figs.~10.6--10.9 & Rewritten with the corrected results; figures
retaken from the corrected run. \\
\S10.7, pp.~192--195; Fig.~10.14 & Rewritten as above; figure retaken. \\
\S10.8, p.~194 & Flow timing 4.1 s. \\
\S10.9, p.~196; \S11.2.1, p.~198 & Corrected numbers; ``from one upload'' becomes ``from one
conversion''. \\
\S11.6, p.~204 & The closing remark now states that the pump which bounds throughput in the
degraded state lies on the restoration critical chain, in place of ``different
components''. \\
Table~B.5, p.~220 & Flow rows corrected. \\
\bottomrule
\end{tabularx}

\end{document}
